\section{Data}\label{sec:Data}

We use proprietary administrative data from Statistics Netherlands. The data on the biannual standardized tests is gathered by the Netherlands Cohort Study on Education (NCO) and subsequently made available to Statistics Netherlands. The NCO collected the test score data only from schools using the tests developped by \textit{Cito}. Since 2018–19, data were obtained by individually contacting primary schools, each of which had to grant permission. As a result, the test score data covers about 50\% of all primary schools in the Netherlands. See \cite{haelermans2020using} for details on the NCO data collection. This test score data can be linked to standard administrative data using an anonymized personal ID, allowing us to also observe teacher track recommendations, parental income and education, and several other student background characteristics.

The test score history, from grade one to six, is available for the students who were enrolled in sixth grade from 2018-19 onward. Our sample includes four cohorts: students who are in sixth grade in the school years 2018-19, 2019-20, 2020-21, and 2021-22. The standard administrative data covers 686,309 students across these four cohorts, with almost no missing values. The biannual test score data includes 318,630 students, although not all students took both standardized tests in each grade. After merging these two datasets, we observe 276,841 students in 5,802 schools. After dropping students with one or both grade 5 test scores missing, as well as those with extreme household parental income values (discussed below), our main estimation sample consists of 148,019 students in 3,295 schools.

\subsection{Descriptive statistics}

\begin{table}
\small
\linespread{1.00}
\begin{center}
\caption{Descriptive statistics}
\label{tbl:descrstats}
\begin{tabular}{lcccc|ccc}
\toprule \toprule
&\multicolumn{4}{c}{Estimation sample}&\multicolumn{3}{c}{Samples w/ missings}\\
[-.75em]
&\multicolumn{4}{c}{}&\multicolumn{3}{c}{Admin \, Test score \, Merged}\\ 
&\multicolumn{1}{c}{(1)}&\multicolumn{1}{c}{(2)}&\multicolumn{1}{c}{(3)}&\multicolumn{1}{c}{(4)}&\multicolumn{1}{c}{(5)}&\multicolumn{1}{c}{(6)}&\multicolumn{1}{c}{(7)}\\
[-.50em]
& Mean & SD  & Min & Max & \multicolumn{3}{c}{Mean} \\
\midrule
&\multicolumn{7}{c}{Panel A: Teacher recommendation} \\
$\mathbbm{1}$(recom. $\geq$ upper voc. tr.)&       0.770&       0.421&       0.000&       1.000 &0.754&&0.758 \\
$\mathbbm{1}$(recom. $\geq$ college tr.)&       0.497&       0.500&       0.000&       1.000 &0.478&&0.489\\
$\mathbbm{1}$(recom. $\geq$ university tr.)&       0.211&       0.408&       0.000&       1.000 &0.197&&0.209\\
[.25em]
&\multicolumn{7}{c}{Panel B: SES measures} \\
SES income          &       6.439&       3.772&      $<$-0.410 &      $>$55.500 &6.394&&6.472\\
$\mathbbm{1}$(SES inc. $\geq$ median)&       0.517&       0.500&       0.000&       1.000 &0.500&&0.506\\
SES years of schooling&      15.975&       2.334&       8.000&      18.000 &16.004&&15.945\\
$\mathbbm{1}$(SES edu. $\geq$ college)&       0.559&       0.497&       0.000&       1.000 &0.618&&0.553\\
[.25em]
&\multicolumn{7}{c}{Panel C: Test scores} \\
Reading midterm grade 5     &     191.618&      26.766&      $<$81.000&     $>$340.000 &&191.614&192.536\\
Reading end term grade 5    &     197.279&      26.983&       $< $84.000&      $> $354.000 &&195.374&196.367\\
Math midterm grade 5     &     254.142&      27.142&       $< $118.000&      $> $379.000 &&253.982&254.564\\
Math end term grade 5    &     261.959&      26.085&       $< $118.000&      $>$394.000 &&262.583&263.179\\
Avg. midterm grade 5     &       0.005&       0.900&      -5.177&       5.077 &&&\\
Avg. end term grade 5    &       0.004&       0.905&      -5.438&       5.130 &&&\\
[.25em]
&\multicolumn{7}{c}{Panel D: Control variables} \\
$\mathbbm{1}$(female)&       0.499&       0.500&       0.000&       1.000 &0.503&&0.503\\
$\mathbbm{1}$(Western imm.)&       0.076&       0.265&       0.000&       1.000 &0.077&&0.077\\
$\mathbbm{1}$(non-Western imm.)&       0.200&       0.400&       0.000&       1.000 &0.178&&0.198\\
Age                 &      11.590&       0.647&        $< $9.250&      $>$14.160 &11.704&&11.660\\
[.25em]
Observations        &    \multicolumn{4}{c}{148019}      &686309 &318630&  276841  \\
\bottomrule \bottomrule
\end{tabular}
\begin{tablenotes}
\footnotesize
\item[] Notes: this table shows the descriptive statistics across four samples. Income is measured in \texteuro 10,000.
%
We translated the schooling levels in the Netherlands to years of schooling as follows: primary school $=$ 8, mbo 1 $=$ 13 , havo or vwo $=$ 14, mbo 2 or 3 $=$ 15, mbo 4 $=$ 16, college or university bachelor $=$ 17, and university master $=$ 18 years.
%
Statistics Netherlands makes a distinction between individuals with a Western and non-Western migration background. Non-Western immigrants have a background from all countries in Africa, Latin America or Asia (including Turkey, excluding Indonesia and Japan), and Western immigrants from the remaining countries.
%
Statistics Netherlands prevents reporting minima and maxima that are based upon less than 10 observations. Therefore, the minimum (maximum) may contain the smallest (largest) number, such that at least ten observations have values lower (higher) than that number. 
\end{tablenotes}
\end{center}
\end{table}


The first four columns in \Cref{tbl:descrstats} present the descriptive statistics for our estimation sample. Panel A focuses on the measures used for the subjective teacher evaluation. Our primary measure is a dummy variable that equals one if the initial track recommendation is at least equal to the college track. College or university track enrollment is highly policy relevant, since it implies the student is tracked towards higher education. Nearly half (49.7\%) of the students receive a track recommendation equal to, or higher than, the college track.

Panel B presents the statistics for the different measures of SES. Our baseline SES measure is gross annual household parental income, which averages \euro64,390. We dropped observations with a yearly parental income above (below) the 99.9th (0.1st) percentile. In our analysis below we standardize parental income per cohort. For robustness checks, we use three alternative SES measures: parental years of schooling of the highest educated parent, a binary income indicator denoting whether parental income exceeds the student's cohort median, and a binary education indicator reflecting whether the highest educated parent holds at least a college degree.

Panel C shows the statistics for the biannual \textit{Cito} tests in fifth grade, per domain and the average across domains. The domain-specific scores are the raw scores estimated on a common scale by the IRT model. The average score is calculated by first standardizing the domain-specific scores per cohort, and subsequently taking the mean across domains. \Cref{tbl:descrstats_tests} in Appendix D provides the descriptive statistics for the raw domain-specific scores on the biannual tests in grades two through four. We exclude grade 1 data since a relatively high fraction of the test scores for that grade are missing. Panel D presents the statistics for the variables used as controls in the robustness checks.

We also compare the statistics of the students in our estimation sample with those from three larger samples representing the different steps of our selection process. Columns (5) to (7) correspond to the samples from the standard administrative data covering the full population, the test score data, and the merged dataset, respectively. The statistics in these columns are based on different numbers of observations due to missing values in some variables, particularly in the grade 5 test scores of columns (6) and (7). The comparison indicates that our estimation sample closely resembles the student population in the Netherlands.